% preamble.tex - shared by main.tex and by every per-chapter test wrapper.
% Engine: Tectonic (XeTeX). Do not add packages in chapter files; add them here.

\usepackage[a4paper,top=25mm,bottom=25mm,inner=28mm,outer=22mm,headheight=14pt]{geometry}
\usepackage{fontspec}
\usepackage{lmodern}
\usepackage{amsmath,amssymb,amsthm,mathtools,bm}
\usepackage{microtype}
\usepackage[dvipsnames,svgnames,table]{xcolor}
\usepackage{graphicx}
\usepackage{booktabs,longtable,array,tabularx,multirow}
\usepackage{enumitem}
\usepackage{caption}
\usepackage{subcaption}
\usepackage{float}
\usepackage{url}
\usepackage{listings}
\usepackage{tikz}
\usetikzlibrary{arrows.meta,positioning,calc,patterns,decorations.pathreplacing,shapes.geometric,fit,backgrounds,matrix,shapes.misc}
\usepackage{pgfplots}
\pgfplotsset{compat=1.18}
\usepgfplotslibrary{fillbetween,groupplots}
\usepackage[most]{tcolorbox}
\usepackage{fancyhdr}
\usepackage{imakeidx}
\makeindex[intoc,columns=2,title=Index]
\usepackage[numbers,sort&compress]{natbib}
\usepackage[hidelinks,bookmarksnumbered]{hyperref}
\usepackage[nameinlink,capitalise,noabbrev]{cleveref}

\graphicspath{{figures/}{figures/svg/}}
\setlist{itemsep=2pt,topsep=4pt}
\setlength{\parskip}{3pt}
\captionsetup{font=small,labelfont=bf}

% ---------------------------------------------------------------- colours
% The course figures' colour language: teal spectral, indigo spatial,
% plum attention, amber loss/error, grey masked.
\definecolor{spec}{HTML}{0F766E}
\definecolor{spat}{HTML}{3F4CA8}
\definecolor{attn}{HTML}{8E3B8A}
\definecolor{loss}{HTML}{B45309}
\definecolor{maskgrey}{HTML}{9CA3AF}
\definecolor{boxone}{HTML}{E6F2F1}
\definecolor{boxwork}{HTML}{F5F3FF}
\definecolor{boxwarn}{HTML}{FEF3C7}
\definecolor{boxsrc}{HTML}{FDECEC}

% ---------------------------------------------------------------- headers
\pagestyle{fancy}
\fancyhf{}
\fancyhead[LE]{\small\leftmark}
\fancyhead[RO]{\small\rightmark}
\fancyfoot[C]{\thepage}
\renewcommand{\headrulewidth}{0.4pt}

% ---------------------------------------------------------------- inline code and paths
% \code{...} and \file{...} accept underscores and break at punctuation.
% They are url-style commands: do NOT use them inside \section, \caption or
% \index arguments - use \texttt{...} with \_ escaped there instead.
\Urlmuskip=0mu plus 2mu
\makeatletter\g@addto@macro{\UrlBreaks}{\UrlOrds\do\_}\makeatother
\DeclareUrlCommand\code{\urlstyle{tt}}
\DeclareUrlCommand\file{\urlstyle{tt}}
\newcommand{\shape}[1]{\texttt{#1}}   % tensor shapes, e.g. \shape{(B, 32, H, W)}

\lstset{basicstyle=\ttfamily\small,columns=fullflexible,keepspaces=true,
  breaklines=true,frame=single,rulecolor=\color{maskgrey},
  backgroundcolor=\color{gray!4},xleftmargin=4pt,xrightmargin=4pt,
  showstringspaces=false,language=Python,literate={~}{{\textasciitilde}}1,
  commentstyle=\color{spec},keywordstyle=\color{spat}\bfseries}

% ---------------------------------------------------------------- maths macros
\newcommand{\vect}[1]{\mathbf{#1}}          % vectors: bold lower case
\newcommand{\mat}[1]{\mathbf{#1}}           % matrices: bold upper case
\newcommand{\R}{\mathbb{R}}
\newcommand{\E}{\mathbb{E}}
\newcommand{\T}{^{\top}}                    % transpose
\newcommand{\inv}{^{-1}}
\newcommand{\norm}[1]{\left\lVert #1 \right\rVert}
\newcommand{\abs}[1]{\left\lvert #1 \right\rvert}
\newcommand{\ip}[2]{\left\langle #1, #2 \right\rangle}
\newcommand{\Cov}{\boldsymbol{\Sigma}}       % covariance matrix
\newcommand{\muv}{\boldsymbol{\mu}}          % mean vector
\newcommand{\xhat}{\hat{\vect{x}}}
\newcommand{\diag}{\operatorname{diag}}
\newcommand{\tr}{\operatorname{tr}}
\newcommand{\softmax}{\operatorname{softmax}}
\newcommand{\argmax}{\operatorname*{arg\,max}}
\newcommand{\argmin}{\operatorname*{arg\,min}}
\newcommand{\sign}{\operatorname{sign}}
\newcommand{\GELU}{\operatorname{GELU}}
\newcommand{\LN}{\operatorname{LN}}
\newcommand{\SAM}{\operatorname{SAM}}

% ---------------------------------------------------------------- theorem-like
\theoremstyle{definition}
\newtheorem{definition}{Definition}[chapter]
\crefname{definition}{Definition}{Definitions}

% ---------------------------------------------------------------- teaching boxes
% 1. The one thing this chapter teaches - first thing in every chapter.
\newtcolorbox{onething}{enhanced,breakable,colback=boxone,colframe=spec,
  fonttitle=\bfseries,title={The one thing this chapter teaches},
  left=6pt,right=6pt,top=4pt,bottom=4pt,before skip=8pt,after skip=10pt}

% 2. Worked example: numbered per chapter, listed in the List of Worked Examples.
%    Usage: \begin{worked}{Mahalanobis distance of a 2-band pixel}\label{we:slug:x} ... \end{worked}
\newtcolorbox[auto counter,number within=chapter,list inside=loe,
  crefname={worked example}{worked examples},
  Crefname={Worked example}{Worked examples}]{worked}[2][]{enhanced,breakable,
  colback=boxwork,colframe=attn!70!black,fonttitle=\bfseries,
  title={Worked example~\thetcbcounter: #2},list entry={#2},
  left=6pt,right=6pt,top=4pt,bottom=4pt,before skip=8pt,after skip=8pt,#1}

% 3. Source-versus-course disagreements (the source wins).
\newtcolorbox{sourcenote}[1][]{enhanced,breakable,colback=boxsrc,colframe=red!55!black,
  fonttitle=\bfseries,title={Source check#1},left=6pt,right=6pt,top=4pt,bottom=4pt}

% 4. Analogy box.
\newtcolorbox{analogy}{enhanced,breakable,colback=white,colframe=spat!60,
  fonttitle=\bfseries,title={Analogy},left=6pt,right=6pt,top=4pt,bottom=4pt,
  boxrule=0.5pt}

% 5. Key point: the one sentence to remember in a section.
\newtcolorbox{keypoint}{enhanced,breakable,colback=boxwarn,colframe=loss,
  left=6pt,right=6pt,top=3pt,bottom=3pt,boxrule=0.6pt,sharp corners}

% 6. Values that will drift (dated numbers from configs, checkpoints, runs).
\newtcolorbox{drift}{enhanced,breakable,colback=gray!6,colframe=maskgrey,
  fonttitle=\bfseries,title={Values that will drift},left=6pt,right=6pt,top=3pt,bottom=3pt}

% 7. Abridged-chapter notice.
\newtcolorbox{abridged}{enhanced,breakable,colback=boxsrc,colframe=red!55!black,
  fonttitle=\bfseries,title={Abridged chapter},left=6pt,right=6pt}

% 8. Lesson box: a mistake the project made, what it cost, and the habit it teaches.
%    Usage: \begin{lesson}[: short title] ... \end{lesson}
\newtcolorbox{lesson}[1][]{enhanced,breakable,colback=boxwarn!60!white,colframe=loss!80!black,
  fonttitle=\bfseries,title={Lesson#1},left=6pt,right=6pt,top=4pt,bottom=4pt}

% Common confusions: \begin{confusions} \confusion{question}{answer} ... \end{confusions}
\newenvironment{confusions}{\section{Common confusions}\begin{description}[style=nextline,leftmargin=1.2em,font=\normalfont\bfseries]}{\end{description}}
\newcommand{\confusion}[2]{\item[#1] #2}

% Check yourself: \begin{checkyourself} \item ... \end{checkyourself}
% Each \item is an exercise numbered c.n; put the solutions in answers/<same-name>.tex.
\newcounter{exercise}[chapter]
\renewcommand{\theexercise}{\thechapter.\arabic{exercise}}
\crefname{exercise}{Exercise}{Exercises}
\newenvironment{checkyourself}{\section{Check yourself}
  \begin{list}{\textbf{Exercise~\theexercise.}}{\usecounter{exercise}\setlength{\leftmargin}{0pt}\setlength{\itemindent}{0pt}\setlength{\labelwidth}{0pt}\setlength{\itemsep}{6pt}}}{\end{list}}

% Answers appendix: each answers file starts with \answersfor{ch:slug}
% then \begin{solution}{ex:slug:1} ... \end{solution}
\newcommand{\answersfor}[1]{\section*{Chapter~\ref{#1}: \nameref{#1}}\addcontentsline{toc}{section}{Chapter~\ref{#1}: \nameref{#1}}}
\newenvironment{solution}[1]{\par\medskip\noindent\textbf{Solution to \cref{#1}.}\quad}{\par}

% Chapter-end pointer.
\newcommand{\nextchapter}[1]{\par\bigskip\noindent\textit{Next: \cref{#1}, \nameref{#1}.}}

% Glossary entries (backmatter): \glossentry{term}{definition}
\newcommand{\glossentry}[2]{\item[#1] #2}

% Figure from the course SVG set: \coursefig[width]{file}{caption}{label}
\newcommand{\coursefig}[4][\linewidth]{\begin{figure}[htbp]\centering
  \includegraphics[width=#1]{#2}\caption{#3}\label{#4}\end{figure}}

% Provenance tags for project numbers:
%   \tagM{research/experiments/x/README.md}  measured, with its source file
%   \tagC                                    computed for this book (by hand in the text, or a script)
%   \tagA                                    assumed for illustration (state the assumption in the text)
\newcommand{\tagM}[1]{\,{\footnotesize\textcolor{gray}{[M:\,\urlstyle{tt}\url{#1}]}}}
\newcommand{\tagC}{\,{\footnotesize\textcolor{gray}{[C]}}}
\newcommand{\tagA}{\,{\footnotesize\textcolor{gray}{[A]}}}
%   \tagMr{text}  measured, but the source is a report not (yet) in a repo file - say which
\newcommand{\tagMr}[1]{\,{\footnotesize\textcolor{gray}{[M:\,#1]}}}

\setcounter{secnumdepth}{2}
\setcounter{tocdepth}{1}
